\documentclass[12pt, a4paper]{article}
\usepackage[utf8]{inputenc}
\usepackage[T1]{fontenc}
\usepackage[english]{babel}
\usepackage{geometry}
\geometry{a4paper, margin=1in}
\usepackage{cite}
\usepackage{booktabs}     % Professional table formatting
\usepackage{siunitx}      % SI units formatting
\usepackage{microtype}    % Typography improvements

\title{\textbf{Compartment Attractor Neural Network (ComAN):} Few-shot learning neural model without backprop for MNIST.}
\author{Mykyta Lapin \and Kostiantyn Bokhan}
\date{\today}

\begin{document}

\maketitle

\begin{abstract}
    Modern deep neural networks (DNNs) deliver impressive accuracy yet depend on backpropagation and large labeled datasets, while offering limited transparency. We present the \emph{Energy–Structural Method} (ESM) for few-shot structural classification, replacing error-driven weight optimization with explicit internal-energy minimization and interpretable attractor formation. Concretely, ESM converts images into graphs via a skeletonization pipeline (binarization $\rightarrow$ skeleton extraction $\rightarrow$ Growing Neural Gas graphing $\rightarrow$ Ramer–Douglas–Peucker simplification), then performs iterative structural–parametric reduction to form a class concept (an energy attractor) under Lyapunov-style operators—eliminating backpropagation and its vanishing-gradient pathologies \cite{rumelhart1986,hochreiter2001}. In preliminary MNIST experiments on six structurally diverse digits (1, 2, 3, 6, 7, 9), using only 5–6 unique examples per class (at most 36 total) augmented to 350 training instances and a single one-pass epoch without backpropagation, ESM attains \emph{Accuracy 82.44\%, Precision 83.33\%, Recall 82.44\%, F1 82.25\%} on 4{,}734 unseen test images. Compared to few-shot/small-data baselines, these results are competitive while using far fewer unique samples and no hyperparameter tuning, in contrast to gradient-based methods that often require extensive data, multiple epochs, and careful regularization \cite{dietterich1995,lipton2018,kaplan2020,crick1989}. Overall, ESM constitutes a viable alternative to classical DNNs for few-shot learning: it is transparent (concepts are explicit hypergraph-like attractors), robust (one-pass learning), and effective from small samples, aligning model dynamics with energy-minimization principles rather than backpropagation.
\end{abstract}

\newpage

\section{Introduction}

Recent advances in artificial intelligence, particularly in deep learning and artificial neural networks (ANNs), have achieved remarkable progress in solving complex problems. However, the widespread deployment of these technologies has revealed fundamental limitations that challenge their ability to create truly autonomous, adaptive, and intelligent systems capable of self-organization and flexible interaction with dynamic environments.

These limitations manifest in several critical areas: the requirement for massive datasets alongside substantial computational and human resources for training and operation \cite{strubell2019, thompson2020}; significant energy consumption and environmental impact \cite{patterson2021}; and the frequent occurrence of errors and ``hallucinations'' in generative models \cite{ji2023, huang2023}, which significantly reduce trust in their outputs and limit practical applications. We argue that these challenges stem from inherent constraints within the existing conceptual framework that are fundamental rather than temporary in nature.

Modern neural network paradigms are primarily grounded in statistical learning approaches that operate on abstract symbolic representations—bits, tokens, and code tables—that serve as substitutes for actual energy transformations. This abstraction creates what we term the ``entropy gap'' phenomenon in generative models \cite{shumailov2023, alemohammad2023}, where there exists a fundamental disparity between statistically frequent events and rare but meaningful ones. This gap reflects a deeper issue: current systems optimize for probable patterns rather than informative ones, leading to predictable but potentially less creative outputs.

\subsection*{The Energy-Information Paradigm}

The core premise of our research is that information itself possesses an inherently energetic nature that can be formalized through the concept of energy landscapes and thermodynamic principles. Unlike traditional approaches that treat information as abstract symbols, our framework conceptualizes information as structured external energy parameters that govern the formation of a system's internal energetic structure—essentially its model of the external world.

This perspective bridges several fundamental theories. Shannon's information entropy can be unified with thermodynamic entropy when information is viewed as structured energy rather than abstract bits. The relationship becomes apparent when considering Landauer's principle \cite{landauer1961}, which establishes that information erasure requires physical energy expenditure. Our theoretical framework extends this connection by proposing that information processing in neural systems should be understood as energy transformation processes governed by the same principles that guide physical systems toward states of minimal internal energy.

\subsection*{Energy Landscapes and Internal Structure Formation}

Central to our approach is the concept of energy landscapes \cite{wales2003}—mathematical descriptions of how systems evolve in state space toward minimal energy configurations. In biological systems, external energy is directly converted into internal energy through sensory processing, accompanied by the extraction of maximal information through both detection of structural elements and measurement of their multiple parameters. For instance, the visual system distinguishes not only linear segments but also their position, orientation, motion direction, length, thickness, and spatial relationships with other elements \cite{hubel1968}.

These measured parameters are projected onto internal quantitative scales and coordinate systems. To enable gradient descent along the energy landscape during internal energy minimization, we must construct hierarchical structures of these scales, enabling transitions from high-energy quantitative representations to lower-energy generalized or qualitative parameters during learning. This process involves scale segmentation down to ordinal scales and generalized coordinate systems, forming the foundation for energy-based structural and parametric reduction operators.

\subsection*{From Detection to Concept Formation}

Our experimental implementation demonstrates this theoretical framework through a multi-stage process that models biological information processing mechanisms. The preprocessing stage simulates sensory system operation through binarization and skeletonization, creating simplified models of sensory and pre-detector neuron functions. The detection stage extracts structural elements (line segments, angular points, intersections) and measures their energy parameters (length, orientation, position), forming the internal energy structure of the system.

The critical concept formation stage models the dendritic tree processing of detector neurons through hypergraph construction and reduction. This stage implements energy-based structural and parametric reduction through gradient descent along the system's energy landscape, guided not by external error minimization but by internal energy minimization principles. The resulting structures—energy attractors—represent stable states that correspond to learned concepts, formed through the system's self-organization rather than external supervision.

\subsection*{Experimental Validation}

To validate these theoretical principles, we developed the Compartment Attractor Neural Network (ComAN) model, which demonstrates energy-based learning on ultra-small training datasets without backpropagation algorithms. The model achieves 82.44\% accuracy on MNIST digit classification using only 5--6 unique examples per class, trained in a single pass. This performance on minimal data suggests that energy-based self-organization may offer fundamental advantages over statistical learning approaches for creating adaptive, efficient artificial intelligence systems.

\section{Background and Related Work}

The proposed Energy--Structural Model (ESM) builds upon and departs from several established directions in pattern recognition and neural processing.

\subsection*{Structural Pattern Classification}

Classical structural approaches represent patterns as symbolic or graph-based descriptions, employing techniques such as syntactic analysis (e.g., grammar-based parsing of structure) and graph matching (e.g., subgraph isomorphism or edit-distance measures) \cite{deraedt2008}. These methods offer \textbf{high interpretability}, since decision rules map directly to recognizable structural features. However, they suffer from \textbf{combinatorial complexity} (graph matching is often NP-hard) and high \textbf{sensitivity to noise}, as minor perturbations in the input structure can lead to large changes in parsing or graph-matching cost.

\subsection*{Skeleton-Based Methods}

Many symbol-recognition and biometric systems extract the medial axis or ``skeleton'' of shapes as a first step \cite{latecki1998, lam1997}. Typical pipelines binarize the input image; apply thinning or distance-transform--driven skeletonization; convert the resulting skeleton into a graph (e.g., via Growing Neural Gas or other graph-extraction algorithms); and finally perform recognition by matching graph features or applying global graph analyzers. These approaches preserve topological structure and can yield compact representations, but they often require multiple heuristic preprocessing stages and may lose discriminative detail due to aggressiveness of skeleton pruning.

\subsection*{Alternatives Without Backpropagation}

A variety of learning paradigms avoid iterative gradient-descent training:

\begin{itemize}
\item \textbf{Extreme Learning Machines (ELM):} Single-hidden-layer feedforward networks with randomly assigned input weights and analytically determined output weights \cite{huang2006}. Training reduces to a linear system solution, offering extremely fast learning but limited capacity to refine internal representations.

\item \textbf{Liquid State / Echo-State Networks:} Recurrent networks with fixed random connectivity; only readout weights are trained \cite{jaeger2001}. They exploit rich ``reservoir'' dynamics for temporal pattern processing, yet require careful parameter tuning to maintain the echo-state property.

\item \textbf{Hopfield-type models:} Associative-memory networks defined by an energy function whose minima store patterns \cite{hopfield1982}. Updates proceed via energy minimization, yielding one-pass convergence to attractors. While biologically inspiring, classical Hopfield nets have limited storage capacity and rely on symmetric weight matrices, and they do not generalize well to high-dimensional sensory inputs.
\end{itemize}

\subsection*{Positioning the Energy--Structural Model (ESM)}

ESM \textbf{integrates structural graph representations} with \textbf{energy-based self-organization} principles \cite{parzhin2025}. Rather than treating energy functions as abstract loss metrics requiring backpropagation, ESM defines \textbf{internal energy landscapes} directly over structural elements (e.g., hypergraphs of skeleton segments). Learning proceeds by \textbf{Lyapunov-driven reduction operators} that minimize internal energy in a single pass, enabling biologically plausible, one-shot formation of attractors without gradient descent. This combination of graph-based pattern encoding and physically grounded energy minimization distinguishes ESM from both purely structural classifiers and conventional energy-based neural models.

\section{ESM Architecture (2 pages)}
\begin{itemize}
\item \textbf{Overall pipeline:} Detailed walkthrough of the diagram in \textit{docs/images/basic\_architecture.png}; trace data flow from the input image to the final decision.
\item \textbf{Section 3.1: Preprocessing:} Skeletonization and initial contour analysis. Goal: an edge-thickness-invariant graph representation.
\item \textbf{Section 3.2: Concept formation:} Secondary analysis, reductions, and iterative composition. Goal: find a stable attractor as a generalized class representation.
\item \textbf{Section 3.3: Classification and decision:} Inference as reduction of the input pattern to an attractor. Compute activations; concept competition via a winner-takes-all (WTA) rule.
\item \textbf{Data-flow table/schematic:} Pixel mask $\rightarrow$ skeleton $\rightarrow$ primary graph $\rightarrow$ reduced graph $\rightarrow$ concept/decision; list key operations at each stage.
\end{itemize}

\section{Practical Implementation of Components (2 pages)}
\begin{itemize}
\item \textbf{4.1 Skeletonization and graph construction:}
\begin{itemize}
\item \textit{Algorithms:} Morphological skeletonization (e.g., Zhang--Suen) with adaptive thresholding; Growing Neural Gas (GNG) for topology learning; Ramer--Douglas--Peucker (RDP) for simplification (as in \textit{README.md}).
\item \textit{Data structures:} Graph representation (e.g., NetworkX); node attributes (point type, coordinates); edge attributes (length, angle).
\item \textit{Complexity:} Complexity as a function of image size.
\end{itemize}
\item \textbf{4.2 Primary contour analysis:}
\begin{itemize}
\item \textit{Parameters:} Enumerate computed parameters (angles, lengths, quadrants, directions). Clarify absolute vs. relative coordinate systems (per \textit{README.md}).
\item \textit{Normalization:} Normalize coordinates for scale- and shift-invariance.
\end{itemize}
\item \textbf{4.3 Secondary analysis and reductions:}
\begin{itemize}
\item \textit{Rules:} Detail reduction rules (IntersectionPoint $\rightarrow$ CornerPoint $\rightarrow$ EndPoint; endpoint removal). Provide visuals as in \textit{README.md}.
\item \textit{Pseudocode:} Include pseudocode for key procedures, e.g., path-finding to delete an endpoint.
\end{itemize}
\item \textbf{4.4 Iterative composition of concepts:}
\begin{itemize}
\item \textit{Training termination:} In the current implementation, each concept finishes training after processing all available training samples; dynamic stabilization criteria are not applied.
\item \textit{Attractor persistence:} Explain parameter updates for a concept (e.g., moving from absolute values to ranges).
\end{itemize}
\end{itemize}

\section{Learning Mechanism and Concept Base Formation (1 page)}
\begin{itemize}
\item \textbf{Training cycle:} Step-by-step:
\begin{enumerate}
\item Use the first class sample as the seed concept.
\item Take the next sample; reduce its graph to the structure of the current concept.
\item Perform structural and parametric generalization (reduction) of the concept based on the new sample.
\item Repeat for all training samples.
\end{enumerate}
\item \textbf{Data organization:} How concepts are stored (e.g., in a graph DB such as Neo4j, as referenced in code, or the filesystem). Structure of a stored concept.
\item \textbf{Start-point selection:} Strategy for selecting a stable start point via clustering of critical points (per \textit{README.md}). Explain why a stable start point is crucial for consistent analysis.
\end{itemize}

\section{Classification Module (1 page)}
\begin{itemize}
\item \textbf{Inference procedure:} Detailed steps for classifying a new image.
\begin{enumerate}
\item Skeletonization and graph construction.
\item For each concept in the base:
\begin{itemize}
\item Early rejection (complexity filter).
\item Reduce the image graph to the concept structure.
\item Compute the activation metric.
\end{itemize}
\item Choose the class via the Winner-Takes-All (WTA) principle.
\end{enumerate}
\item \textbf{Activation metric:} Specify similarity computation. Use the idea of Graph Edit Distance (GED) as an ``activation energy'' proxy. Formula: $\mathrm{Similarity} = 1 - (\mathrm{GED} / \mathrm{MaxPossibleCost})$.
\item \textbf{Class selection logic:} Describe the WTA mechanism and tie-breaking by concept complexity.
\end{itemize}

\section{Experimental Results (1 page)}
\begin{itemize}
\item \textbf{Setup:} Dataset (MNIST), number of classes, and number of training samples per scenario (no augmentation, +2, +5, +10).
\item \textbf{Detailed tables:} Present tables from \textit{README.md} with metrics (Accuracy, Precision, Recall, F1-score) per scenario.
\item \textbf{Trend analysis:} How metrics change with more training data; emphasize positive dynamics.
\item \textbf{Plot:} Include and comment on \textit{docs/images/augmentation\_impact\_chart.png}, visualizing augmentation impact.
\item \textbf{Discussion:} Why augmentation is important here (it helps form a more robust and generalized attractor).
\end{itemize}

\section{Discussion and Limitations (0.5 page)}
\begin{itemize}
\item \textbf{Strengths:}
\begin{itemize}
\item \textbf{Transparency:} Ability to visualize and analyze formed concepts.
\item \textbf{Forward-only learning:} Computationally efficient training without backpropagation.
\item \textbf{Learning from small data:} Ability to form concepts from few examples.
\end{itemize}
\item \textbf{Application scenarios:} Where the approach is useful (medical diagnostics, handwritten symbol recognition---anywhere structure and interpretability matter).
\item \textbf{Limitations:}
\begin{itemize}
\item \textbf{Scaling:} Behavior with hundreds or thousands of classes.
\item \textbf{Skeletonization quality:} Strong dependence on the first stage; errors there cannot be corrected later.
\item \textbf{Complex textures:} Poor performance where class identity is texture- rather than structure-driven (e.g., animal breed by fur).
\end{itemize}
\item \textbf{Comparison to conventional NNs:} Summarize where ESM wins and where CNNs still dominate.
\end{itemize}

\section{Conclusions and Future Work (0.5 page)}
\begin{itemize}
\item \textbf{Summary:} We designed and validated an architecture for forward-only structural learning.
\item \textbf{Achieved autonomy:} The system forms concepts autonomously, though start-point selection and reduction rules still require manual tuning.
\item \textbf{Roadmap:}
\begin{itemize}
\item Expand to more classes and harder datasets (e.g., EMNIST, Quick, Draw!).
\item Optimize computational complexity of reductions and graph matching.
\item Automate hyperparameter selection (e.g., thresholds for start-point clustering).
\item Develop ``neuron maps'' to model intra-class variability. For instance, the generalized concept ``digit 7'' could combine sub-concepts (attractors) for different handwriting styles (with and without a crossbar).
\end{itemize}
\end{itemize}

\section{References}

\bibliographystyle{IEEEtran}
\begin{thebibliography}{99}

\bibitem{rumelhart1986}
D.~E. Rumelhart, G.~E. Hinton, and R.~J. Williams, ``Learning representations by back-propagating errors,'' \emph{Nature}, vol.~323, pp. 533--536, Oct. 1986.

\bibitem{hochreiter2001}
S.~Hochreiter, Y.~Bengio, P.~Frasconi, and J.~Schmidhuber, ``Gradient flow in recurrent nets: The difficulty of learning long-term dependencies,'' in S.~C. Kremer and J.~F. Kolen (eds.), \emph{A Field Guide to Dynamical Recurrent Neural Networks}. Piscataway, NJ, USA: IEEE Press, 2001, pp. 237--243.

\bibitem{lipton2018}
Z.~C. Lipton, ``The mythos of model interpretability,'' \emph{Queue}, vol.~16, no.~3, pp. 31--57, Jun. 2018.

\bibitem{dietterich1995}
T.~G. Dietterich, ``Overfitting and undercomputing in machine learning,'' \emph{ACM Computing Surveys}, vol.~27, no.~3, pp. 326--327, Sep. 1995.

\bibitem{kaplan2020}
J.~Kaplan, S.~McCandlish, T.~Henighan, T.~B. Brown, B.~Chess, R.~Child, S.~Gray, A.~Radford, J.~Wu, and D.~Amodei, ``Scaling laws for neural language models,'' \emph{arXiv preprint} arXiv:2001.08361, 2020.

\bibitem{crick1989}
F.~Crick, ``The recent excitement about neural networks,'' \emph{Nature}, vol.~337, pp. 129--132, Jan. 1989.

\bibitem{strubell2019}
E.~Strubell, A.~Ganesh, and A.~McCallum, ``Energy and policy considerations for deep learning in NLP,'' in \emph{Proceedings of the 57th Annual Meeting of the Association for Computational Linguistics}, 2019, pp. 3645--3650.

\bibitem{thompson2020}
S.~Thompson, K.~Greenewald, K.~Lee, and G.~F. Manso, ``The computational limits of deep learning,'' \emph{arXiv preprint} arXiv:2007.05558, 2020.

\bibitem{patterson2021}
D.~Patterson, J.~Gonzalez, Q.~Le, C.~Liang, L.~M. Munguia, D.~Rothchild, D.~So, M.~Texier, and J.~Dean, ``Carbon emissions and large neural network training,'' \emph{arXiv preprint} arXiv:2104.10350, 2021.

\bibitem{ji2023}
Z.~Ji, N.~Lee, R.~Frieske, T.~Yu, D.~Su, Y.~Xu, E.~Ishii, Y.~J. Bang, A.~Madotto, and P.~Fung, ``Survey of hallucination in natural language generation,'' \emph{ACM Computing Surveys}, vol.~55, no.~12, pp. 1--38, 2023.

\bibitem{huang2023}
L.~Huang, W.~Yu, W.~Ma, W.~Zhong, Z.~Feng, H.~Wang, Q.~Chen, W.~Peng, X.~Feng, B.~Qin, and T.~Liu, ``A survey on hallucination in large language models: Principles, taxonomy, challenges, and open questions,'' \emph{arXiv preprint} arXiv:2311.05232, 2023.

\bibitem{shumailov2023}
I.~Shumailov, Z.~Shumaylov, Y.~Zhao, Y.~Gal, N.~Papernot, and R.~Anderson, ``The curse of recursion: Training on generated data makes models forget,'' \emph{arXiv preprint} arXiv:2305.17493, 2023.

\bibitem{alemohammad2023}
S.~Alemohammad, J.~Casco-Rodriguez, L.~Luzi, A.~I. Humayun, H.~Babaei, D.~LeJeune, A.~Siahkoohi, and R.~G. Baraniuk, ``Self-consuming generative models go MAD,'' \emph{arXiv preprint} arXiv:2307.01850, 2023.

\bibitem{landauer1961}
R.~Landauer, ``Irreversibility and heat generation in the computing process,'' \emph{IBM Journal of Research and Development}, vol.~5, no.~3, pp. 183--191, 1961.

\bibitem{wales2003}
D.~J. Wales, \emph{Energy Landscapes}. Cambridge, UK: Cambridge University Press, 2003.

\bibitem{hubel1968}
D.~H. Hubel and T.~N. Wiesel, ``Receptive fields and functional architecture of monkey striate cortex,'' \emph{The Journal of Physiology}, vol.~195, no.~1, pp. 215--243, 1968.

\bibitem{deraedt2008}
L.~De Raedt, \emph{Logical and Relational Learning}. Berlin, Germany: Springer, 2008.

\bibitem{latecki1998}
L.~J. Latecki and R.~Lakaemper, ``Shape similarity measure based on correspondence of visual parts,'' \emph{IEEE Trans. Pattern Anal. Mach. Intell.}, vol.~22, no.~10, pp. 1185--1190, 2000.

\bibitem{lam1997}
L.~Lam and C.~Y.~Suen, ``An evaluation of parallel thinning algorithms for character recognition,'' \emph{IEEE Trans. Pattern Anal. Mach. Intell.}, vol.~17, no.~9, pp. 914--919, 1995.

\bibitem{huang2006}
G.-B. Huang, Q.-Y. Zhu, and C.-K. Siew, ``Extreme learning machine: Theory and applications,'' \emph{Neurocomputing}, vol.~70, no.~1--3, pp. 489--501, Dec. 2006.

\bibitem{jaeger2001}
H.~Jaeger, ``The `echo state' approach to analysing and training recurrent neural networks,'' GMD Report 148, German National Research Center for Information Technology, 2001.

\bibitem{hopfield1982}
J.~J. Hopfield, ``Neural networks and physical systems with emergent collective computational abilities,'' \emph{Proceedings of the National Academy of Sciences}, vol.~79, no.~8, pp. 2554--2558, Apr. 1982.

\bibitem{parzhin2025}
Y.~Parzhin, M.~Lapin, and K.~Bokhan, ``A new approach to building energy models of neural networks,'' \emph{Energy-Based Neural Network Modeling}, preprint, 2025.

\end{thebibliography}


\section{Appendices (up to 1 page)}
\begin{itemize}
\item Provide detailed schemata of reduction rules.
\item Show 2--3 example graphs for MNIST digits before and after full reduction to a concept.
\item Include concise pseudocode for the most complex algorithms (e.g., iterative concept update).
\end{itemize}

\end{document}
